\documentclass[11pt]{article}
\usepackage[T1]{fontenc}
\usepackage[utf8]{inputenc}
\usepackage{lmodern}
\usepackage{amsmath,amssymb}
\usepackage{longtable,booktabs,array}
% Compatibility with Pandoc's captionless tables on older longtable versions.
\ifcsname c@none\endcsname\else\newcounter{none}\fi
\usepackage{graphicx}
\usepackage[margin=25mm]{geometry}
\usepackage{hyperref}
\hypersetup{colorlinks=true,linkcolor=blue,urlcolor=blue}
\providecommand{\tightlist}{\setlength{\itemsep}{0pt}\setlength{\parskip}{0pt}}
\providecommand{\pandocbounded}[1]{#1}
\setlength{\emergencystretch}{3em}
\setcounter{secnumdepth}{-1}
\begin{document}
\section{Serre's comparison theorems and Chow's
theorem}\label{serres-comparison-theorems-and-chows-theorem}

\emph{Written and self-checked by GPT-6.1 Sol (OpenAI) in Codex, Ultra
setting, October 2026. Analytic prerequisite integration by GPT-6 Astra
(OpenAI), Codex, Ultra. The linked analytic bridge was written and
self-checked by Claude Opus 5.5; a full independent review is not
claimed. Independently authored material: CC0.}

Locally, analytification preserves exact sequences but adds holomorphic
functions. On a projective variety, the three comparison theorems say
much more: coherent cohomology agrees, every analytic morphism of
coherent sheaves is algebraic, and every coherent analytic sheaf is
algebraic. We prove these assertions in that order. The proof of the
last assertion needs a finite-dimensional analytic cohomology theorem;
using the comparison theorem to assert finiteness for an arbitrary
analytic sheaf at this stage would be circular.

Throughout, varieties are reduced separated schemes of finite type over
\(\mathbf C\), as in
\href{complex-analytic-spaces-and-analytification.html}{Complex analytic
spaces and analytification}. Write \(P=\mathbf P^n_{\mathbf C}\),
\(P^{\mathrm{an}}=\mathbf P^n(\mathbf C)\), and
\(\mathcal O(d)^{\mathrm{an}}\) for the analytification of a twist. We
use the generation, vanishing and projective cohomology proved in
\href{serres-theorems-on-projective-schemes.html}{Serre's theorems} and
\href{cohomology-of-projective-space.html}{Cohomology of projective
space}.

\subsection{1. The analytic cohomology
foundation}\label{the-analytic-cohomology-foundation}

Two analytic theorems are required for the cohomology arguments below.
Their programme proofs are in \emph{Complex analytic spaces and coherent
sheaves}. Read the vanishing and finiteness lessons before this lesson:
neither obtains analytic cohomology from GAGA or from algebraization.
Jean-Pierre Demailly's freely available
\href{https://www-fourier.univ-grenoble-alpes.fr/~demailly/manuscripts/agbook.pdf}{\emph{Complex
Analytic and Differential Geometry}} is further reading. The
\href{prerequisites.html}{prerequisite and proof guide} identifies the
exact statements, proof locations and remaining dependencies.

\textbf{Analytic vanishing theorem.} If a complex manifold admits a
smooth strictly plurisubharmonic exhaustion, every coherent analytic
sheaf on it has zero cohomology in positive degrees. An exhaustion has
relatively compact sublevel sets; strict plurisubharmonicity means its
complex Hessian is positive definite.
\href{../complex-analytic-spaces-and-coherent-sheaves/theorems-a-and-b-on-stein-manifolds.html\#4-theorem-b}{Theorems
A and B on Stein manifolds, Theorem 4.1}, supplies the proof through
local coherent resolutions, approximation and exhaustion. The assertion
applies to every coherent sheaf, not only vector bundles; it gives
vanishing in every positive degree, not merely finite dimension.

\textbf{Analytic finiteness theorem (Cartan--Serre).} On a compact
complex analytic space, the cohomology of every coherent analytic sheaf
is finite-dimensional in every degree.
\href{../complex-analytic-spaces-and-coherent-sheaves/finiteness-on-compact-complex-spaces.html\#2-the-finiteness-theorem}{Finiteness
on compact complex spaces, Theorem 2.1}, gives the proof using finite
nested acyclic covers. Restriction of coherent sections from a larger
open set to a relatively compact smaller one is a compact operator, by
the coherent version of Montel's theorem. The induced map between their
Čech complexes is an isomorphism on cohomology. The compact-operator
criterion is developed in the preceding lesson,
\href{../complex-analytic-spaces-and-coherent-sheaves/frechet-spaces-of-sections-and-schwartzs-theorem.html}{Fréchet
spaces of sections and Schwartz's theorem}. It yields finite dimension
and closed coboundaries. These arguments apply to coherent analytic
sheaves before any algebraization is known.

For our standard projective cover \(U_i=\{T_i\ne0\}\), every nonempty
intersection is biholomorphic to

\[
(\mathbf C^*)^r\times\mathbf C^{n-r}.
\]

The function

\[
\psi(z)=\sum_{j=1}^n|z_j|^2+\sum_{j=1}^r|z_j|^{-2}
\]

is a strictly plurisubharmonic exhaustion there. The first sum makes the
Hessian positive, while the second prevents escape toward a deleted
coordinate hyperplane. The vanishing theorem therefore makes this cover
acyclic for every coherent analytic sheaf. By the Čech comparison proved
earlier in the course, its ordered complex computes cohomology. There
are \(n+1\) opens, so

\[
H^q(P^{\mathrm{an}},M)=0\qquad(q>n)
\tag{1}
\]

for every coherent analytic module \(M\). This bound concerns all
coherent analytic modules, including those not yet known to be
algebraic.

\subsection{2. Starting the comparison: the structure sheaf and
twists}\label{starting-the-comparison-the-structure-sheaf-and-twists}

\textbf{Lemma 2.1.} The analytic structure sheaf of projective space has
cohomology \(\mathbf C\) in degree zero and zero in every positive
degree. The natural comparison from algebraic cohomology is an
isomorphism.

\textbf{Proof.} For \(n=0\), the space is a point. For general \(n\),
use the acyclic cover of Section 1. Put
\(\Omega_S=\{T\in\mathbf C^{n+1}:T_i\ne0\text{ for }i\in S\}\). A
function on \(U_S=\bigcap_{i\in S}U_i\) pulls back to a homogeneous
holomorphic function of weight zero on this product of punctured planes
and planes. Conversely a weight-zero function descends by setting any
\(T_j=1\), \(j\in S\); homogeneity makes this independent of the
representative. Work directly on \(\Omega_S\), without choosing
projective ratios for the expansion.
\href{../../human/elementary-analysis/laurent-series-and-homogeneous-projections.html\#2-products-of-annuli-and-discs}{Laurent
series and homogeneous projections, Theorems 1--4}, proves the unique,
normally convergent product expansion and its homogeneous decomposition.
The allowed exponents are

\[
\alpha\in\mathbf Z^{n+1},\qquad
\sum_i\alpha_i=0,\qquad
\alpha_j\geq0\quad(j\notin S).
\]

Indeed uniqueness applied to \(f(2T)=f(T)\) makes the coefficient of
\(T^\alpha\) vanish unless its total exponent is zero. Group the
expansion by \(N=\{j:\alpha_j<0\}\). There are only finitely many such
sets.
\href{../../human/elementary-analysis/laurent-series-and-homogeneous-projections.html\#3-extending-the-parts-with-specified-negative-powers}{Theorem
3 of the Laurent unit} proves that each part extends to the whole
\(\Omega_N\): on any compact subset, choose smaller positive contour
radii in the negative-power coordinates and larger contour radii in the
others. The product geometric estimate proves normal convergence there,
including along the newly allowed zero coordinates. The extension still
has weight zero. Uniqueness of coefficients on \((\mathbf C^*)^{n+1}\)
makes every such projection compatible with restriction between chart
intersections. Thus each fixed \(N\) has one coefficient space,
independent of the containing face \(S\), rather than a separate choice
for every chart.

For fixed \(N\), the allowed faces of the Čech complex are exactly the
nonempty \(S\) containing \(N\). If \(N\) is proper and nonempty, choose
\(j\notin N\). Insert \(j\) with the sign of its ordered position,
extending the component on \(\Omega_{S\cup\{j\}}\) first to
\(\Omega_N\), then restricting to \(\Omega_S\). If \(j\) is already in
the target face, set the contraction to zero.
\href{../../human/elementary-analysis/laurent-series-and-homogeneous-projections.html\#4-homogeneous-functions-and-the-ordered-complex}{Theorem
5} gives the signed formula and proves \(dh+hd=1\), including degree
zero. The term removing the inserted \(j\) is the identity; every other
pair differs by one in its insertion/removal sign and cancels. This is
an operation on normally convergent holomorphic components, not only a
formal monomial calculation.

If \(N\) is empty, all exponents are nonnegative and sum to zero, so the
only monomial is 1. This is the ordinary simplex complex of constants,
with cohomology \(\mathbf C\) in degree zero. If \(N\) were all indices,
their sum would be negative, so this part is absent. These contractions
prove the analytic assertion. Algebraically the same cohomology was
computed in our projective-space lesson, and comparison sends a constant
to the identical constant. \(\square\)

For any algebraic module \(F\), the map on global sections to those of
\(F^{\mathrm{an}}\) extends canonically to maps

\[
c_F^q:H^q(X,F)\longrightarrow H^q(X^{\mathrm{an}},F^{\mathrm{an}}).
\tag{2}
\]

Indeed analytification is an exact functor on all modules, by local
flatness. The target groups therefore form a cohomological delta
functor, and the universal derived functors of algebraic global sections
extend the degree-zero map uniquely. Consequently (2) is natural and
commutes with connecting maps in long exact sequences. For coherent
modules, it also commutes with closed-immersion pushforward, whose
stalkwise quotient description was proved in the preceding lesson.

\textbf{Lemma 2.2.} Comparison is an isomorphism for \(\mathcal O(d)\)
on \(\mathbf P^n\), for every integer \(d\) and every degree.

\textbf{Proof.} Induct on \(n\). On \(\mathbf P^0\) every twist is a
one-dimensional module, including negative twists. For \(n>0\), let
\(i:E\hookrightarrow P\) be a hyperplane. Its equation gives

\[
0\to\mathcal O(d-1)\to\mathcal O(d)
\to i_*\mathcal O_E(d)\to0.
\tag{3}
\]

Analytification preserves this sequence. The induction hypothesis
supplies comparison on \(E\). The two long exact sequences show that
comparison for \(\mathcal O(d)\) in all degrees is equivalent to
comparison for \(\mathcal O(d-1)\): for either unknown term, take five
successive terms with that term in the middle; all four other comparison
maps are isomorphisms. Begin with \(d=0\), proved in Lemma 2.1, and move
upward and downward. \(\square\)

\subsection{3. The first comparison
theorem}\label{the-first-comparison-theorem}

\textbf{Theorem 3.1 (cohomological GAGA).} If \(X\) is projective over
\(\mathbf C\) and \(F\) is coherent, (2) is an isomorphism for every
\(q\geq0\).

\textbf{Proof.} Embed \(X\) as a closed subvariety of \(P\).
Closed-immersion pushforward is exact, preserves coherence and
cohomology, and commutes with analytification. Thus it suffices to prove
the result on \(P\).

Descend on \(q\), simultaneously for all coherent sheaves. For \(q>n\),
both groups vanish: algebraically by the standard affine cover,
analytically by (1). Choose a presentation

\[
0\to R\to L\to F\to0,
\tag{4}
\]

where \(L\) is a finite sum of twists and \(R\) is coherent. Suppose
comparison in degree \(q+1\) is already known for every coherent sheaf.
Lemma 2.2 gives comparison for \(L\) in every degree.

To prove surjectivity for \(F\) in degree \(q\), start with an analytic
class. Its boundary in \(H^{q+1}(R^{\mathrm{an}})\) lifts uniquely to an
algebraic class. That lift maps to zero in \(H^{q+1}(L)\), since its
analytic image does and comparison for \(L\) is injective. Exactness
lifts it to \(H^q(F)\). The difference between this class's analytic
image and the original class has zero boundary, so comes from
\(H^q(L^{\mathrm{an}})\). Comparison for \(L\) lifts the difference,
establishing surjectivity.

This proves surjectivity in degree \(q\) for every coherent sheaf, in
particular for \(R\). If an algebraic class in \(H^q(F)\) has zero
analytic image, its boundary is zero by injectivity in degree \(q+1\)
for \(R\), so it comes from \(H^q(L)\). That lift maps analytically into
the image of \(H^q(R^{\mathrm{an}})\). Surjectivity for \(R\) lifts this
correcting class. Subtract it in \(H^q(L)\); comparison for \(L\) forces
the corrected lift to be zero. Thus the original class was zero. The
induction proves injectivity and surjectivity in all degrees.
\(\square\)

The simultaneous induction is essential: we first establish surjectivity
for all coherent sheaves before using it for the kernel \(R\). No finite
global locally free resolution of \(F\) has been assumed.

\subsection{4. The second comparison
theorem}\label{the-second-comparison-theorem}

\textbf{Theorem 4.1 (full faithfulness).} For coherent \(F,G\) on a
projective complex variety,

\[
\operatorname{Hom}_X(F,G)
\xrightarrow{\sim}
\operatorname{Hom}_{X^{\mathrm{an}}}(F^{\mathrm{an}},G^{\mathrm{an}}).
\]

\textbf{Proof.} The internal Hom \(A=\mathcal Hom(F,G)\) is coherent.
Finite presentations and local flatness give

\[
A^{\mathrm{an}}\cong\mathcal Hom(F^{\mathrm{an}},G^{\mathrm{an}}),
\]

as proved in Proposition 5.1 of the preceding lesson. Apply Theorem 3.1
in degree zero to \(A\). These global sections are exactly the two
morphism groups, and the comparison takes a morphism to its
analytification. \(\square\)

In particular, an analytic isomorphism between algebraic coherent
sheaves lifts uniquely to an algebraic isomorphism: lift it and its
inverse, then use faithfulness to identify their compositions with the
identities.

\subsection{5. The third comparison
theorem}\label{the-third-comparison-theorem}

\textbf{Lemma 5.1 (analytic generation).} Every coherent analytic module
\(M\) on \(P^{\mathrm{an}}\) has \(M(d)\) generated by finitely many
global sections for all sufficiently large \(d\).

\textbf{Proof.} We prove this together with essential surjectivity,
inducting on \(n\). At \(n=0\) a coherent module is a finite-dimensional
vector space. Assume essential surjectivity on every hyperplane
\(E\cong\mathbf P^{n-1}\). On \(E\), Theorem 3.1 and algebraic Serre
vanishing now imply vanishing of all positive cohomology of sufficiently
large twists of every coherent analytic module.

Fix \(x\in P^{\mathrm{an}}\) and choose a hyperplane through \(x\), with
equation \(t\). Multiplication has coherent kernel and cokernel:

\[
0\to C\to M(-1)\xrightarrow{t}M\to B\to0.
\]

Both \(C\) and \(B\) are annihilated by \(t\), so are coherent modules
on \(E^{\mathrm{an}}\). This remains true when multiplication by \(t\)
is not injective. Let \(L_d\) be the image of \(M(d-1)\to M(d)\). The
two short exact sequences give, for large \(d\), surjections

\[
H^1(M(d-1))\twoheadrightarrow H^1(L_d)
\twoheadrightarrow H^1(M(d)),
\tag{5}
\]

because \(H^2(C(d))=H^1(B(d))=0\). By the Cartan--Serre finiteness
theorem in Section 1, these dimensions are finite. Thus the nonnegative
integers \(h^1(M(d))\) eventually stop decreasing. In the stable range
both arrows in (5) are isomorphisms. The second short exact sequence
then shows that

\[
H^0(M(d))\twoheadrightarrow H^0(B(d)).
\tag{6}
\]

By induction \(B\) is algebraic on \(E\); its large twists are generated
by global sections. Lift generators using (6). At \(x\), these generate
\(M(d)_x/tM(d)_x\); Nakayama makes them generate \(M(d)_x\), since \(t\)
belongs to the local maximal ideal.

Generation by finitely many fixed sections holds on a neighbourhood,
because their coherent cokernel vanishes there. Once generation holds at
a point in degree \(d\), multiplying these sections by a homogeneous
coordinate nonzero at that point proves generation in every degree
\(d'\geq d\). Compactness supplies finitely many neighbourhoods and a
common bound on \(d\). At each such degree take the union of their
finite generating families. This proves the lemma. \(\square\)

\textbf{Theorem 5.2 (essential surjectivity).} Every coherent analytic
module on \(X^{\mathrm{an}}\), for projective \(X\), is isomorphic to
\(F^{\mathrm{an}}\) for some coherent algebraic \(F\). Together with
Theorem 4.1, analytification is an equivalence of coherent categories.

\textbf{Proof.} First take \(X=P\), continuing the induction of Lemma
5.1. Global generation gives a surjection \(L_0^{\mathrm{an}}\to M\),
where \(L_0\) is a finite sum of equal negative twists. Its kernel \(R\)
is coherent. Generate a twist of \(R\) as well, to obtain

\[
L_1^{\mathrm{an}}\xrightarrow{v}L_0^{\mathrm{an}}\to M\to0.
\]

Full faithfulness lifts \(v\) uniquely to an algebraic map
\(u:L_1\to L_0\). Put \(F=\operatorname{coker}u\). Exactness of
analytification gives \(F^{\mathrm{an}}\cong M\), completing the
induction on \(n\).

For \(i:X\hookrightarrow P\), algebraize \(i_*M\) to \(G\) on \(P\). If
\(I\) is the ideal of \(X\), then
\((IG)^{\mathrm{an}}=I^{\mathrm{an}}G^{\mathrm{an}}=0\). Exactness and
faithfulness give \(IG=0\). Hence \(G=i_*F\) for a unique coherent
module \(F\) on \(X\). Compatibility with closed-immersion pushforward
identifies \(F^{\mathrm{an}}\) with \(M\). \(\square\)

The correct uniqueness assertion fixes the identification. If
\((F,\phi)\) and \((G,\psi)\) are algebraizations equipped with
isomorphisms to \(M\), there is a unique isomorphism \(u:F\to G\)
satisfying \(\psi\circ u^{\mathrm{an}}=\phi\). Without specified
identifications, an algebraic sheaf can have many automorphisms: every
nonzero scalar acts on \(\mathcal O_X\). Thus its bare isomorphism class
is unique, but its possible isomorphisms are not.

\subsection{6. Chow's theorem and geometric
consequences}\label{chows-theorem-and-geometric-consequences}

\textbf{Theorem 6.1 (Chow).} Every closed analytic subset of a complex
projective variety is the analytification of a unique reduced closed
algebraic subvariety.

\textbf{Proof.} Its analytic vanishing ideal
\(J\subset\mathcal O_{X^{\mathrm{an}}}\) is coherent by Cartan
coherence. Algebraize \(J\) as a coherent module by Theorem 5.2, and
lift its inclusion into the structure sheaf by Theorem 4.1. Exactness
and faithfulness make the lift injective, with image an algebraic
coherent ideal \(I\) whose analytification is \(J\).

The ideal \(I\) is radical. At every closed stalk, if \(a^m\in I_x\),
the image of \(a\) belongs to the radical ideal
\(J_x=I_x\mathcal O_{X^{\mathrm{an}},x}\); faithful flatness contracts
this to \(a\in I_x\). Radicality at closed stalks suffices: a nonzero
nilpotent ideal in the coherent quotient would have support containing a
closed point. The reduced closed scheme defined by \(I\) therefore
analytifies to the given analytic subset. If two algebraic ideals have
the same analytic extension, faithful flatness gives equal closed
stalks, and coherent support detects equality globally. This proves
uniqueness. \(\square\)

\textbf{Theorem 6.2.} Every holomorphic map between the analytifications
of projective complex varieties is the analytification of a unique
algebraic morphism.

\textbf{Proof.} Its graph is a closed reduced analytic subspace of the
projective product. Apply Chow's theorem to algebraize the graph to
\(Z\subset X\times Y\). The first projection \(p:Z\to X\) is proper and
its analytification is an isomorphism. Each closed fiber has one complex
point and is zero-dimensional. Upper semicontinuity of fiber dimension
for proper maps shows that all fibers are zero-dimensional: any nonempty
closed locus of positive-dimensional fibers would contain a closed
point. Thus \(p\) has finite fibers and is finite by
\href{the-theorem-on-formal-functions.html\#5-dimension-and-finiteness-consequences}{Formal
functions, Theorem 5.2}.

At a closed point \(x\), the finite algebra \((p_*\mathcal O_Z)_x=B\)
over \(A=\mathcal O_{X,x}\) has only one maximal ideal, since the closed
fiber has one point. Its maximal-ideal topology is cofinal with its
\(\mathfrak m_A\)-adic topology. Thus
\(B\otimes_A\widehat A=\widehat B\). Local analytic comparison and the
analytic graph isomorphism identify the map \(\widehat A\to\widehat B\)
with an isomorphism. Faithful flatness of completion makes \(A\to B\) an
isomorphism. The kernel and cokernel of
\(\mathcal O_X\to p_*\mathcal O_Z\) are coherent and vanish at every
closed point, so they vanish everywhere. A finite map is the relative
spectrum of this algebra, hence \(p\) is an isomorphism. The second
projection composed with \(p^{-1}\) is the required morphism. Uniqueness
follows from uniqueness of the reduced graph. \(\square\)

\textbf{Proposition 6.3.} For \(n\geq1\), every analytic line bundle on
\(\mathbf P^n(\mathbf C)\) is \(\mathcal O(d)^{\mathrm{an}}\) for a
unique \(d\in\mathbf Z\).

\textbf{Proof.} GAGA gives a coherent algebraic module \(L\). It is a
line bundle: at a closed point choose one element lifting an analytic
fiber basis. Nakayama and faithful flatness make the resulting map
\(A\to L_x\) onto; after analytic tensor it is a map between rank-one
free modules with nonzero residue, hence an isomorphism. Faithfulness
makes its algebraic kernel zero. The locally free locus is open, and its
closed complement contains no closed point, so is empty.

A rational nonzero section of \(L\) gives a Cartier divisor. Projective
space is locally factorial, since its affine local rings are
localizations of polynomial rings. Every prime divisor is the zero set
of an irreducible homogeneous polynomial \(f\): its homogeneous prime
cone ideal has height one in the factorial ring
\(\mathbf C[T_0,\ldots,T_n]\), so is principal. If \(\deg f=e\) and
\(H=V(T_0)\), then the rational function \(f/T_0^e\) has divisor
\(V(f)-eH\). Summing shows every divisor is linearly equivalent to
\(dH\), and hence \(L\cong\mathcal O(d)\). If \(\mathcal O(d)\) were
trivial, its global-section dimension would be one. The calculation in
our projective-space lesson gives zero for \(d<0\), and
\(\binom{d+n}{n}\) for \(d\geq0\), which is one only for \(d=0\). This
proves uniqueness. On \(\mathbf P^0\), there is just the trivial line
bundle and every twist represents it. \(\square\)

A compact Riemann surface holomorphically embedded in projective space
consequently has algebraic image. It is an algebraic smooth projective
curve: dimensions and regularity agree at complex points by local
comparison. The hypothesis includes an embedding; this conclusion does
not assert that an arbitrary compact analytic space is projectively
embeddable.

\subsection{7. Exercises with solutions}\label{exercises-with-solutions}

\textbf{Exercise 7.1 (easy).} Show why the cohomological comparison
theorem fails for \(\mathbf A^1\).

\textbf{Solution.} In degree zero it is the inclusion
\(\mathbf C[z]\to\mathcal O(\mathbf C)\). The entire function \(e^z\) is
absent from the image, as its derivatives never become zero. Thus even
degree-zero surjectivity fails. The projective hypothesis supplies the
compactness and twist arguments used above.

\textbf{Exercise 7.2 (medium).} Classify analytic line bundles on
projective space, explaining the dimension-zero exception.

\textbf{Solution.} Essential surjectivity algebraizes the bundle; local
faithful flatness descends its rank-one freeness as in Proposition 6.3.
A rational section and homogeneous prime-divisor equations make its
divisor a multiple of a hyperplane, so it is
\(\mathcal O(d)^{\mathrm{an}}\). For \(n\geq1\), the global-section
dimension distinguishes the trivial twist and hence makes \(d\) unique.
For \(n=0\), every \(\mathcal O(d)\) is the same one-dimensional bundle,
so there is no unique integer.

\textbf{Exercise 7.3 (medium).} Deduce Chow's theorem for a nonempty
analytic hypersurface of \(\mathbf P^n\) directly from its ideal, using
GAGA.

\textbf{Solution.} An analytic hypersurface in projective space has an
invertible coherent ideal: the convergent local rings are factorial, by
\href{../complex-analytic-spaces-and-coherent-sheaves/the-local-ring-of-holomorphic-germs.html\#3-unique-factorization}{The
local ring of holomorphic germs, Theorem 3.1}, so a reduced pure
codimension-one vanishing ideal is locally the product of its prime
equations. Algebraize this ideal and its inclusion into the structure
sheaf. The ideal descends to an algebraic line bundle by the same local
argument as in Proposition 6.3, so is \(\mathcal O(-d)\). Its inclusion
is a nonzero section of \(\mathcal O(d)\), hence a homogeneous
polynomial of degree \(d\). A nonempty zero set requires \(d>0\). The
ideal's analytic extension is the original vanishing ideal; radicality
descends by faithful flatness. Thus this polynomial defines exactly the
given reduced hypersurface; the ideal formulation handles reducible
hypersurfaces as well.

\textbf{Exercise 7.4 (medium).} Describe holomorphic maps
\(\mathbf P^1(\mathbf C)\to\mathbf P^1(\mathbf C)\).

\textbf{Solution.} Theorem 6.2 makes the map algebraic. Its pullback of
\(\mathcal O(1)\) is \(\mathcal O(d)\), and the pulled-back coordinate
sections generate it, so \(d\geq0\). They are homogeneous degree-\(d\)
polynomials \(F_0,F_1\) without a common projective zero. The map is
\([T_0:T_1]\mapsto[F_0(T):F_1(T)]\), and in an affine target chart it is
their rational quotient. For \(d=0\) it is constant; for \(d>0\) its
degree is \(d\).

\textbf{Exercise 7.5 (medium).} Compute analytic cohomology of
\(\mathcal O(d)\) on \(\mathbf P^1\), and verify comparison explicitly.

\textbf{Solution.} Use \(z=T_1/T_0\), with trivializations
\(e_0=T_0^d\), \(e_1=T_1^d=z^de_0\). The acyclic two-chart cover gives
the degree-one quotient of holomorphic functions on \(\mathbf C^*\) by
entire functions in \(z\) and \(z^d\) times entire functions in
\(z^{-1}\). In the Laurent expansion the first family removes exponents
\(a\geq0\); the second removes exponents \(a\leq d\). The remaining
basis is \(z^{d+1},\ldots,z^{-1}\) when \(d\leq-2\), and is empty
otherwise. Thus \(h^1=\max(-d-1,0)\). Global sections must have
exponents both \(a\geq0\) and \(a\leq d\), giving \(1,z,\ldots,z^d\) for
\(d\geq0\) and zero otherwise. Therefore \(h^0=\max(d+1,0)\); higher
cohomology vanishes. The algebraic Laurent Čech calculation has the
identical bases, and comparison takes each basis monomial to itself.

\textbf{Exercise 7.6 (harder).} State and prove the uniqueness of
algebraization with a fixed analytic identification. Give a
counterexample to uniqueness of the bare isomorphism.

\textbf{Solution.} For \(\phi:F^{\mathrm{an}}\to M\) and
\(\psi:G^{\mathrm{an}}\to M\), full faithfulness lifts \(\psi^{-1}\phi\)
uniquely to \(u:F\to G\). Lifting the inverse and using faithfulness
proves \(u\) is an isomorphism. It is the unique one compatible with the
identifications. On a nonempty projective variety, multiplication by any
\(\lambda\in\mathbf C^*\) is an automorphism of \(\mathcal O_X\); there
are many isomorphisms of that bare sheaf to itself. The compatibility
condition singles out the correct one.

\subsection{8. Sources and the comparison
chain}\label{sources-and-the-comparison-chain}

J.-P. Serre's
\href{https://www.numdam.org/item/AIF_1956__6__1_0/}{\emph{Géométrie
algébrique et géométrie analytique}}, Annales de l'Institut Fourier 6
(1956), 1--42, is the historical source for these three theorems and
their applications. The linked edition is free to read; no permission to
copy or translate it is asserted. Chow's theorem is credited to W.-L.
Chow. The arguments here are independently written; Section 1 links the
analytic programme proofs and Demailly's freely accessible exposition.
Component authorship and reuse terms remain recorded in the respective
lessons.

Kiran S. Kedlaya's
\href{https://ocw.mit.edu/courses/18-726-algebraic-geometry-spring-2009/resources/mit18_726s09_lec22_gaga/}{\emph{GAGA},
updated 30 April 2009}, MIT OpenCourseWare 18.726, treats comparison,
full faithfulness, algebraization and Chow in Sections 3--5 and 8. Its
notes are reusable under
\href{https://ocw.mit.edu/pages/privacy-and-terms-of-use/}{CC BY-NC-SA
4.0}. They offer useful parallel reading, while this course retains its
explicit Laurent contraction, two-stage descending comparison induction
and treatment of noninjective hyperplane multiplication.

For a broader framework, Jack Hall's
\href{https://arxiv.org/abs/1804.01976v3}{\emph{GAGA theorems},
arXiv:1804.01976v3, 17 May 2022}, Theorem A and Examples 9.2--9.3,
compares analytic and formal GAGA for proper schemes. Theorem A requires
coherence, finite total coherent cohomology, detection at closed points
and specified local flatness and residue-field conditions. Its
non-Noetherian extensions have further hypotheses. This public author
draft is further reading, not a source licensed for adaptation or a
replacement for the proofs above.

The corresponding verified open AI Integrated Stacks Project treatments
are
\href{https://kokunoyumeto.github.io/stacks-zh-hans-cn/en/gaga.html\#gaga-section-proof-gaga-cohomology}{cohomological
comparison},
\href{https://kokunoyumeto.github.io/stacks-zh-hans-cn/en/gaga.html\#gaga-section-proof-gaga-fully-faithful}{full
faithfulness},
\href{https://kokunoyumeto.github.io/stacks-zh-hans-cn/en/gaga.html\#gaga-lemma-projective-analytic-global-generation}{analytic
global generation},
\href{https://kokunoyumeto.github.io/stacks-zh-hans-cn/en/gaga.html\#gaga-section-completion-gaga-essential-surjectivity}{essential
surjectivity}, and
\href{https://kokunoyumeto.github.io/stacks-zh-hans-cn/en/gaga.html\#gaga-theorem-chow-closed-analytic-projective}{Chow's
theorem}. Their reference texts retain GFDL 1.2 and contain AI-written
additions not reviewed by the official Stacks maintainers. Section 1
states the analytic finiteness and acyclicity inputs; Section 2 gives
the Laurent structure-sheaf calculation assuming those inputs. The
referenced chapter also invokes analytic finiteness in its generation
argument.

Using the analytic programme results in Section 1 and the preceding
lesson, the arguments above establish comparison and algebraicity for
projective reduced complex varieties. The chain is local analytic
algebra and faithful flatness; analytic Čech cohomology and finiteness;
comparison for twists; comparison for coherent sheaves; internal Hom and
full faithfulness; analytic generation and two presentations;
algebraicity of ideals and graphs. GAGA for nonreduced projective
schemes requires the nonreduced analytification construction. GAGA for
arbitrary proper complex schemes additionally requires a full proper
reduction and descent argument. These extensions are explicit remaining
proof obligations, rather than consequences proved by the
further-reading references.
\end{document}
